% !PRINTED 595
\weberpage{595}
%<<CONT
from which it follows that \gesperrt{in one of the systems} (subsidiary groups)
\[ X,\ X\psi_0,\ X\psi_0^2,\ \ldots \]
\gesperrt{every substitution $\psi$ of $\Psi$ is contained.}

If $\mathfrak{p}$ is of degree $f$ (with respect to $R$), then
\[ \omega^{P^f} \equiv \omega\ (\mathrm{mod}\ \mathfrak{p}), \]
and $P^f$ is the lowest power of $P$ with positive exponent that satisfies this congruence for all $\omega$.

If we therefore set $\nu = f$ in (10), it follows that $\psi_0^f$ is contained in $X$, but that no power of $\psi_0$ with a lower positive exponent has this property. The subsidiary groups
\[ X,\ X\psi_0,\ X\psi_0^2,\ \ldots,\ X\psi_0^{f-1} \]
are therefore all different from one another, and we obtain:
\weq{13}{\Psi = X + X\psi_0 + X\psi_0^2 + \cdots + X\psi_0^{f-1}.}

By (12), however, the totality of the substitutions $\psi_0^\nu X$ is identical with $X\psi_0^\nu$ (for every $\nu$); hence
\weq{14}{\psi_0^\nu X = X\psi_0^\nu,}
from which it follows that \gesperrt{$X$ is a normal divisor of $\Psi$.}

\gesperrt{Since, as we have seen above, $\Psi$ is of degree $gf$, $X$ is of degree $g$.}

Now in the group $X$ there are substitutions $\chi_1$, among them always the identical one, for which the congruence
\weq{15}{\tau\mkern1mu|\mkern1mu\chi_1 \equiv \tau\ (\mathrm{mod}\ \mathfrak{p}^2)}
is satisfied, and these form a group $X^{(1)}$, whose degree may be denoted by $g_1$.

\gesperrt{It is to be proved that $g_1$ is a power of $p$.} For if we denote by $\pi$ a prime functional divisible by $\mathfrak{p}$, and by $\alpha$ an integral functional, then it follows from (15):
\weq{16}{\tau\mkern1mu|\mkern1mu\chi_1 = \tau + \pi^2\alpha,}
and since likewise, by (15),
\[ \alpha\mkern1mu|\mkern1mu\chi_1 \equiv \alpha\ (\mathrm{mod}\ \mathfrak{p}^2),\quad (\pi\mkern1mu|\mkern1mu\chi_1)^2 \equiv \pi^2\ (\mathrm{mod}\ \mathfrak{p}^3) \]
holds, it results from (16) that
\[ \tau\mkern1mu|\mkern1mu\chi_1^2 \equiv \tau + 2\pi^2\alpha\ (\mathrm{mod}\ \mathfrak{p}^3), \]
and in general, for every positive exponent $h$,
\[ \tau\mkern1mu|\mkern1mu\chi_1^h \equiv \tau + h\pi^2\alpha\ (\mathrm{mod}\ \mathfrak{p}^3). \]
